% TEMPLATE-MILC-v3.tex - plantilla maestra UNICA de las guias MILC v3.
%
% La usan las cuatro familias de guias, cada una con su driver:
%   · Grados 6-11          -> scripts/build-guias-g11.py
%   · Bebras               -> scripts/build-guias-bebras.py
%   · Semillero            -> scripts/build-guias-semillero.py
%   · Territorio interior  -> scripts/build-guias-territorio-interior.py
%
% Antes habia tres copias casi identicas (~400 lineas iguales, 6 distintas):
% cada mejora habia que aplicarla tres veces, y en la practica no se hacia
% (el motor de recursos vivia solo en dos de ellas). Lo que varia entre
% familias esta parametrizado en 6 placeholders que cada driver llena
% (se nombran aqui SIN su sintaxis real: el builder reemplaza por texto plano,
% tambien dentro de los comentarios, y un valor multilinea rompe el preambulo):
%
%   HEADER_IZQ        encabezado izquierdo de pagina
%   HEADER_DER        encabezado derecho de pagina
%   PORTADA_ETIQUETA  rotulo sobre el numero grande (GRADO/MOMENTO/MODULO)
%   PORTADA_SERIE     linea de serie de la portada
%   PORTADA_ID        identificador de la guia en la portada
%   PORTADA_CIRCULO   el circulito de grado (°), o vacio si no aplica
%   PDF_TITULO        titulo en texto plano (sin \textbf ni comillas TeX) para
%                     los metadatos del PDF (/Title); lo produce lib_xelatex.texto_pdf
%
% Composicion (auditoria 2026-09-03): las bandas de estacion piden espacio
% (\Needspace*) y prohiben el salto justo despues (\nopagebreak), asi no quedan
% huerfanas al pie; las cajas largas (softbox, drawbox, infoband, puentebox) son
% `breakable`, asi una caja de media pagina ya no arrastra todo a la siguiente
% dejando la anterior al 40 %. Todo texto sobre mostaza o verde va en milcNegro
% (blanco sobre #E5B400 da 1,93:1 y sobre #58A923 2,95:1; ninguno pasa WCAG AA).
% El resto de claves las documenta content/guias/_SCHEMA.md. El script valida
% que no queden placeholders sin reemplazar antes de escribir el archivo final.

% Etiquetado PDF (latex-lab): /Lang, estructura y texto alternativo de las
% figuras, para lectores de pantalla. Debe ir ANTES de \documentclass.
\DocumentMetadata{lang=es-CO,tagging=on}
\documentclass[11pt,letterpaper]{article}
% headheight=24pt: la cabecera derecha lleva dos lineas (identidad de la guia
% y estacion actual). headsep se acorta en la misma medida para que el bloque
% de texto quede exactamente donde estaba con headheight=12pt.
\usepackage[margin=2.0cm,headheight=24pt,headsep=13pt]{geometry}
\usepackage[table]{xcolor}
\usepackage{fontspec}
\usepackage[spanish]{babel}
\usepackage{tikz}
% Librerías TikZ para los diagramas de las guías (bloque `recursos:`):
%  · babel  → babel en español vuelve activos caracteres como «>», y TikZ los usa
%             en sus opciones (>=stealth, ->). Sin esta librería, cualquier
%             diagrama con flechas revienta con «\language@active@arg> has an
%             extra }». Debe ir ANTES de que se dibuje nada.
%  · positioning        → sintaxis «below left=2pt and 3pt».
%  · decorations.pathreplacing → llaves ({brace}) para señalar rangos.
\usetikzlibrary{babel,positioning,decorations.pathreplacing}
\usepackage{graphicx}
% Circuitos: el trabajo de electrónica se hace hoy con micro:bit y sus sensores
% integrados (MakeCode, sin cableado), así que no se necesitan esquemáticos.
% Si algún día entran montajes con protoboard, basta descomentar esta línea:
% los diagramas .tex/.tikz declarados en `recursos:` ya se incrustan solos.
% \usepackage{circuitikz}
\usepackage[most]{tcolorbox}
\usepackage{tabularx}
\usepackage{array}
\usepackage{fancyhdr}
\usepackage{titlesec}
\usepackage[document]{ragged2e}  % bandera a la derecha en todo el cuerpo
\usepackage{enumitem}
\usepackage{microtype}
\usepackage{etoolbox}
\usepackage{needspace}
% hyperref al final del preambulo: metadatos del PDF (titulo, autor, idioma,
% familia), marcadores por estacion (\pdfbookmark) y enlaces sin recuadro.
\usepackage[hidelinks,
  pdftitle={Historia de los medios — del pregonero al chat de WhatsApp},
  pdfauthor={PhD. Álvaro Cárdenas Orozco},
  pdfsubject={Tecnología e Informática · Grado 6 · Guía 5 · MILC v3},
  pdflang={es-CO},
  bookmarksopen=true,bookmarksnumbered=false]{hyperref}

% Helvetica Neue no trae la flecha U+2192, asi que cada «→» escrito literalmente
% en el YAML se compilaba a NADA, en silencio (462 flechas en 61 guias: rutas del
% tipo «paleta → tipografia → lockup» salian rotas). Mapeamos el caracter a su
% version matematica, que si tiene glifo. Asi el autor puede seguir escribiendo
% «→» comodamente en el YAML.
\usepackage{newunicodechar}
\newunicodechar{→}{\ensuremath{\rightarrow}}
% Mismo problema con los simbolos de marca y vinetas: el «✓» de «una palomita ✓»
% salia como cuadrito vacio en el PDF de todas las guias que lo usaban.
\usepackage{pifont}
\newunicodechar{✓}{\ding{51}}
\newunicodechar{✗}{\ding{55}}
\newunicodechar{✔}{\ding{52}}
\newunicodechar{•}{\textbullet}
\newunicodechar{≥}{\ensuremath{\geq}}
\newunicodechar{≤}{\ensuremath{\leq}}

\setmainfont{Helvetica Neue}
\setsansfont{Helvetica Neue}
\setlength{\parindent}{0pt}
\setlength{\parskip}{6pt}
% Sin particion de palabras: en 6.º-7.º una palabra cortada por guion al final
% de linea («Comuni-cacion») frena la lectura mucho mas que una linea corta.
\hyphenpenalty=10000
\exhyphenpenalty=10000
\pretolerance=2000
\tolerance=3000
\setlength{\emergencystretch}{3em}
\linespread{1.10}
\renewcommand{\arraystretch}{1.25}
\setlist{leftmargin=5mm,itemsep=1mm,topsep=1mm}
\newcolumntype{Y}{>{\RaggedRight\arraybackslash}X}

\definecolor{milcMagenta}{HTML}{E6007E}
\definecolor{milcVino}{HTML}{5A0038}
\definecolor{milcTurquesa}{HTML}{0093A5}
\definecolor{milcVerde}{HTML}{58A923}
\definecolor{milcMostaza}{HTML}{E5B400}
\definecolor{milcOcre}{HTML}{B86600}
\definecolor{milcNegro}{HTML}{241816}
\definecolor{milcGris}{HTML}{F7F7F4}
\definecolor{milcLinea}{HTML}{E8E2DD}
\definecolor{milcGradePrimary}{HTML}{0066FF}
\definecolor{milcGradeDark}{HTML}{003D99}
\definecolor{milcGradeSoft}{HTML}{D6E8FF}
\definecolor{milcGradeAccent}{HTML}{CFE7FF}
\definecolor{milcGradeText}{HTML}{FFFFFF}
\color{milcNegro}

\pagestyle{fancy}
\fancyhf{}
% Cabecera de dos lineas: identidad de la guia arriba y, a la derecha y en
% gris, la estacion en curso (\markright desde \milcestacion). Antes de la
% primera estacion la segunda linea va vacia; el \strut mantiene la altura
% para que la linea de arriba no baile entre paginas.
\lhead{\small Tecnología e Informática · Grado 6\\[-1pt]{\footnotesize\strut}}
\rhead{\small Guía 5 · MILC v3\\[-1pt]{\footnotesize\strut\color{milcNegro!65}\rightmark}}
\cfoot{\small\thepage}
\renewcommand{\headrulewidth}{0pt}

% Sobre mostaza (#E5B400) y verde (#58A923) el texto va en milcNegro; sobre el
% resto de la paleta, en blanco. Se usa dentro de fonttitle/fontupper, donde un
% \color posterior manda sobre coltitle.
\newcommand{\milctintasobre}[1]{%
  \ifboolexpr{test{\ifstrequal{#1}{milcMostaza}} or test{\ifstrequal{#1}{milcVerde}}}%
    {\color{milcNegro}}{\color{white}}}

\newtcolorbox{titlebox}[1]{
  enhanced,colback=#1,colframe=#1,arc=7mm,boxrule=0pt,
  left=7mm,right=7mm,top=3mm,bottom=3mm,
  before skip=8pt,after skip=8pt,
  fontupper=\sffamily\bfseries\Large\color{white}
}

\newtcolorbox{softbox}[3]{
  enhanced,breakable,pad at break=3mm,
  colback=#2,colframe=#1,borderline west={4pt}{0pt}{#1},
  arc=2mm,boxrule=.4pt,left=4mm,right=4mm,top=3mm,bottom=3mm,
  before skip=6pt,after skip=8pt,title={#3},coltitle=white,
  colbacktitle=#1,fonttitle=\sffamily\bfseries\milctintasobre{#1},fontupper=\RaggedRight,
  attach boxed title to top left={xshift=3mm,yshift=-2mm},
  boxed title style={arc=2mm, boxrule=0pt}
}

\newtcolorbox{drawbox}[2]{
  enhanced,breakable,pad at break=4mm,
  colback=white,colframe=#1,arc=4mm,boxrule=1pt,
  left=5mm,right=5mm,top=4mm,bottom=4mm,before skip=8pt,after skip=8pt,
  title={#2},coltitle=white,colbacktitle=#1,fonttitle=\sffamily\bfseries\milctintasobre{#1},
  fontupper=\RaggedRight,
  attach boxed title to top left={xshift=4mm,yshift=-2mm},
  boxed title style={arc=3mm, boxrule=0pt}
}

\newtcolorbox{infoband}[1]{
  enhanced,breakable,pad at break=2mm,colback=#1!8,colframe=#1!50,arc=2mm,boxrule=.5pt,
  left=4mm,right=4mm,top=2mm,bottom=2mm,before skip=4pt,after skip=4pt,
  fontupper=\small\RaggedRight
}

% Puente narrativo entre secciones (contrato editorial Conéctate).
% Da continuidad: "Acabas de X. Ahora vas a Y porque Z."
% Visualmente: banda izquierda en color del grado, texto en itálica pequeña.
\newtcolorbox{puentebox}{
  enhanced,breakable,pad at break=2mm,colback=milcGradeSoft,colframe=milcGradeDark,
  borderline west={3pt}{0pt}{milcGradeDark},
  arc=0mm,boxrule=0pt,left=5mm,right=4mm,top=2.5mm,bottom=2.5mm,
  before skip=4pt,after skip=8pt,
  fontupper=\itshape\small\color{milcNegro}\RaggedRight
}
% La flecha va en modo matematico a proposito: Helvetica Neue NO tiene el glifo
% U+2192, asi que \textrightarrow se compilaba a NADA («Missing character: There
% is no → in font Helvetica Neue») y el puente salia sin flecha. En modo
% matematico la flecha viene de la fuente matematica y si se dibuja.
\newcommand{\puente}[1]{%
  \begin{puentebox}%
    \textcolor{milcGradeDark}{$\rightarrow$}\hspace{2mm}#1%
  \end{puentebox}%
}

\newcommand{\linead}[1]{\noindent\color{milcLinea}\rule{#1}{0.5pt}\color{milcNegro}}
\newcommand{\respuesta}[1]{\par\vspace{2mm}\foreach \n in {1,...,#1}{\noindent\color{milcLinea}\rule{\linewidth}{0.5pt}\par\vspace{5mm}}\color{milcNegro}}
\newcommand{\checkbox}{\textcolor{milcGradeDark}{\fbox{\phantom{x}}}\hspace{5pt}}

% ─── Listas aireadas para las guías socioemocionales ────────────────────
% Componer las verificaciones como «\checkbox texto\\» deja el texto que salta
% de línea debajo del cuadrito y apelmaza el bloque. Estos entornos alinean la
% continuación y airean los ítems. Los usa Territorio Interior; cualquier
% familia puede hacerlo.
\newlist{checklista}{itemize}{1}
\setlist[checklista]{
  label=\textcolor{milcGradeDark}{\fbox{\phantom{x}}},
  leftmargin=9mm, labelsep=3.2mm, itemsep=2.4mm,
  topsep=2.5mm, parsep=0pt, partopsep=0pt, before=\RaggedRight
}
\newlist{pasoslista}{enumerate}{1}
\setlist[pasoslista]{
  label=\textcolor{milcGradeDark}{\sffamily\bfseries\arabic*.},
  leftmargin=8mm, labelsep=3mm, itemsep=2.8mm,
  topsep=1mm, parsep=0pt, partopsep=0pt, before=\RaggedRight
}

\newcommand{\phasecard}[4]{
\begin{tcolorbox}[enhanced,colback=#3,colframe=#2,arc=3mm,boxrule=.5pt,height=3.05cm,valign=center,halign=center,left=1.5mm,right=1.5mm,top=2mm,bottom=2mm]
{\sffamily\bfseries\fontsize{22}{24}\selectfont\textcolor{#2}{#1}}\par
{\sffamily\bfseries\small #4}
\end{tcolorbox}
}

% ── Piezas del rediseno 2026 (legibilidad 6.º-11.º) ───────────────────
% Banda de estacion: reemplaza el titlebox macizo. Titulo a la izquierda,
% dato practico (tiempo/modalidad) a la derecha, en una sola linea.
% Nunca huerfana: pide 9 lineas libres (o salta de pagina rellenando con
% \vfill) y prohibe el salto justo despues de la banda. Nueve y no seis:
% la caja partible que sigue necesita ~80pt (titulo adosado, relleno y dos
% lineas) para arrancar en la misma pagina; con menos, tcolorbox la manda
% entera a la siguiente y la banda se queda sola. El penalty 10000 va
% en `after`, ANTES del espacio posterior: un `after skip` (aunque sea 0pt)
% es un \vskip tras la caja, y ese glue es un punto de corte legal que un
% \nopagebreak escrito despues de \end{tcolorbox} ya no cierra. Cada banda
% deja marcador PDF y actualiza la cabecera derecha.
\newcounter{milcestacion}
\newcommand{\milcestacion}[2]{%
\Needspace*{9\baselineskip}%
\stepcounter{milcestacion}%
\pdfbookmark[1]{#2}{milcestacion\themilcestacion}%
\markright{#2}%
\begin{tcolorbox}[enhanced,colback=#1,colframe=#1,arc=2mm,boxrule=0pt,
  left=5mm,right=5mm,top=2.6mm,bottom=2.6mm,before skip=11pt,
  after={\par\nopagebreak\vskip7pt}]
{\sffamily\bfseries\fontsize{15}{18}\selectfont\milctintasobre{#1}#2}
\end{tcolorbox}}

% Tarjeta de paso numerada: sustituye las tablas «Pilar / Aplicacion», que
% eran formato de consulta docente, no de estudiante. El numero va en un
% circulo sobre el borde para que la secuencia se lea de un vistazo.
\newcommand{\milcpaso}[3]{%
\Needspace{4\baselineskip}%
\begin{tcolorbox}[enhanced,colback=white,colframe=milcLinea,arc=1.5mm,boxrule=.9pt,
  left=14mm,right=4mm,top=3mm,bottom=3mm,before skip=4pt,after skip=4pt,
  fontupper=\RaggedRight,
  overlay={\node[anchor=north west,circle,fill=#2,text=white,inner sep=0pt,
    minimum size=7.4mm,font=\sffamily\bfseries\small]
    at ([xshift=3.6mm,yshift=-3.2mm]frame.north west) {#1};}]
#3
\end{tcolorbox}}

% Casilla marcable de tamano util con lapiz (la anterior era \fbox{\phantom{x}},
% de alto variable segun el texto que la seguia).
\newcommand{\milcmarca}{\framebox[3.8mm]{\rule{0pt}{3.4mm}}}

% Fila marcable de lista suelta (checks de la estacion 1).
\newcommand{\milccheckrow}[1]{%
\par\vspace{1.8mm}\noindent
\makebox[6.4mm][l]{\raisebox{-.55mm}{\textcolor{milcNegro}{\milcmarca}}}%
\parbox[t]{\dimexpr\linewidth-6.4mm\relax}{\RaggedRight #1}\par}

% Celda marcable dentro de la rubrica (tabla de autoevaluacion). Misma
% sangria francesa, pero pensada para vivir dentro de una columna X.
\newcommand{\milcopcion}[1]{%
\makebox[5.2mm][l]{\raisebox{-.55mm}{\textcolor{milcNegro}{\milcmarca}}}%
\parbox[t]{\dimexpr\linewidth-5.2mm\relax}{\RaggedRight #1}}

% ── Recursos visuales de la guía ──────────────────────────────────────
% Declarados en el bloque `recursos:` del YAML y emitidos por el builder.
% \guiaFigura   → imagen rasterizada o PDF (foto, carta celeste, montaje).
% \guiaDiagrama → fuente TikZ/circuitikz (.tex) incrustada como vector:
%                 es la vía para circuitos y esquemáticos editables.
% [#1] = texto alternativo (alt) · #2 = ruta relativa al .tex · #3 = pie
% El alt llega al PDF etiquetado (/Alt) y lo lee el lector de pantalla.
\newcommand{\guiaFigura}[3][]{%
\begin{center}
\includegraphics[width=0.88\linewidth,height=0.38\textheight,keepaspectratio,alt={#1}]{#2}\\[1.5mm]
{\footnotesize\itshape\textcolor{milcNegro!75}{#3}}
\end{center}
}

\newcommand{\guiaDiagrama}[2]{%
\begin{center}
\input{#1}\\[1.5mm]
{\footnotesize\itshape\textcolor{milcNegro!75}{#2}}
\end{center}
}

\begin{document}
\thispagestyle{empty}

% ── Portada ───────────────────────────────────────────────────────────
% Antes: un rectangulo de color a pagina completa. Se veia bien en pantalla,
% pero en la fotocopiadora del colegio se comia el toner y salia gris sucio.
% Ahora la banda ocupa el tercio superior y el resto va en blanco.
\begin{tikzpicture}[remember picture,overlay]
  \path[fill=milcGradePrimary]
    (current page.north west) rectangle ([yshift=-10.6cm]current page.north east);

  \node[anchor=north west,text=milcGradeText] at ([xshift=2.0cm,yshift=-1.55cm]current page.north west)
    {\sffamily\bfseries\fontsize{12}{14}\selectfont GRADO};
  \node[anchor=north west,text=milcGradeText,opacity=.85] at ([xshift=2.0cm,yshift=-2.35cm]current page.north west)
    {\sffamily\bfseries\fontsize{8.5}{10}\selectfont SERIE GUÍAS · TECNOLOGÍA E INFORMÁTICA};
  \node[anchor=north east,text=milcGradeText,opacity=.85,align=right] at ([xshift=-2.0cm,yshift=-1.55cm]current page.north east)
    {\sffamily\bfseries\fontsize{9}{12}\selectfont I.E. Sor María Juliana\\Cartago, Valle del Cauca};

  \node[anchor=west,text=milcGradeText] (gradonum) at ([xshift=1.85cm,yshift=-7.1cm]current page.north west)
    {\sffamily\bfseries\fontsize{112}{112}\selectfont 6};
  % El circulito de grado (°) solo aplica a las guias de grado («11°»); en
  % Bebras y Semillero el numero grande es un momento/modulo, no un grado.
  % Se posiciona relativo al nodo `gradonum`, no en coordenadas absolutas.
  \draw[milcGradeText,line width=1.6pt]
    ([xshift=2.0mm,yshift=-4.5mm]gradonum.north east) circle (.15cm);

  \node[anchor=west,text=milcGradeText,text width=11.4cm,align=left] at ([xshift=1.5cm,yshift=0.5cm]gradonum.east)
    {
     {\sffamily\bfseries\fontsize{9}{11}\selectfont\MakeUppercase{Período 1 · Comunicación e identidad digital}}\\[3mm]
     {\sffamily\bfseries\fontsize{21}{25}\selectfont\setlength{\baselineskip}{27pt}Historia de\\[4pt]los medios ---\\[4pt]del pregonero\\[4pt]al WhatsApp}\\[3.5mm]
     {\sffamily\bfseries\fontsize{15}{18}\selectfont Guía 5-6-TIC}
    };
\end{tikzpicture}

\vspace*{9.4cm}
\vfill

{\sffamily\bfseries\fontsize{8}{10}\selectfont\textcolor{milcGradeDark}{LO QUE TE LLEVAS HOY}}

\begin{tcolorbox}[enhanced,colback=white,colframe=milcNegro,arc=2mm,boxrule=1.6pt,
  left=5mm,right=5mm,top=4mm,bottom=4mm,before skip=3pt,after skip=12pt,
  fontupper=\RaggedRight\fontsize{11}{15.5}\selectfont]
Una \textbf{línea de tiempo dibujada en hoja completa} del cuaderno, con \textbf{8 medios de comunicación} en orden cronológico (1850 hasta hoy). Cada medio lleva: año aproximado, dibujo pequeño, 1 frase de qué hacía, y \textbf{una flecha que conecta dos medios diciendo qué se conservó} (saludo, asunto, firma, audiencia, etc.). Más \textbf{una conclusión escrita} de 3-4 líneas: ¿qué cambió y qué sigue igual desde el pregonero hasta WhatsApp?
\end{tcolorbox}

\begin{tcolorbox}[enhanced,colback=milcGris,colframe=milcGris,arc=2mm,boxrule=0pt,
  left=5mm,right=5mm,top=3.5mm,bottom=3.5mm,after skip=12pt]
{\sffamily\bfseries\fontsize{8}{10}\selectfont\textcolor{milcNegro!55}{ESTA GUÍA ES DE}}\\[3mm]
\begin{tabularx}{\linewidth}{X p{3.2cm} p{3.2cm}}
\linead{\linewidth} & \linead{3.0cm} & \linead{3.0cm}\\[-1mm]
{\fontsize{8.5}{10}\selectfont\textcolor{milcNegro!55}{Nombre completo}} &
{\fontsize{8.5}{10}\selectfont\textcolor{milcNegro!55}{Grupo}} &
{\fontsize{8.5}{10}\selectfont\textcolor{milcNegro!55}{Fecha}}\\
\end{tabularx}
\end{tcolorbox}

\vfill

\noindent\textcolor{milcLinea}{\rule{\linewidth}{1pt}}\\[2mm]
{\sffamily\bfseries\fontsize{9.5}{12}\selectfont Autor: Dr. Álvaro Cárdenas Orozco}\\[1mm]
{\sffamily\fontsize{8.5}{11}\selectfont\textcolor{milcNegro!55}{Metodología MILC · apertura ancestral · pensamiento computacional · triángulo Dussel–estoicismo–Floridi}}

\newpage

\section*{Apertura: Historia de los medios --- del pregonero al chat de WhatsApp}

\begin{softbox}{milcGradeDark}{milcGradeSoft}{Saber ancestral}
\textbf{Tu bisabuela escuchaba radio --- y antes de la radio escuchaba al pregonero.} Si pudieras viajar al Cartago de 1920, no había radio en cada casa. Había \textbf{el pregonero} en el centro, había \textbf{el periódico} (\emph{El Diario del Valle}) que pasaba el cartero los lunes, y había \textbf{las cartas familiares} que viajaban semanas. Después, en los años 40, llegó la radio: la gente se reunía alrededor del aparato a oír noticias y radionovelas. En los años 60 entró la televisión en blanco y negro, y las familias del barrio se juntaban en la casa del único vecino que la tenía. En los 90 llegó internet a los café-internet. En el 2010 explotó WhatsApp. \textbf{Cada paso de esa historia trajo algo nuevo, pero también guardó algo del paso anterior.} La radio guardó la voz del pregonero. WhatsApp guarda el saludo de la carta. Internet no nació de la nada: es el último capítulo de una historia muy larga.
\end{softbox}

\begin{softbox}{milcTurquesa}{milcGris}{Saber contemporáneo}
Un \textbf{medio de comunicación} es cualquier herramienta que las personas usan para mandarse mensajes a distancia o a muchas personas a la vez. Hoy hay muchos --- WhatsApp, TikTok, YouTube, Instagram --- pero \textbf{todos vienen de medios más antiguos}. Saber esa historia te ayuda a entender 2 cosas: (1) las cosas que parecen \emph{``modernas''} muchas veces vienen de hace siglos --- el asunto del correo viene de la carta postal; (2) cada medio nuevo trae \emph{problemas nuevos} pero también \emph{cuidados antiguos} que ya conocemos. Hoy vas a recorrer 8 medios desde 1850 hasta hoy y vas a descubrir el hilo que los une.
\end{softbox}

\begin{softbox}{milcMostaza}{milcGris}{Pregunta-puente}
\textit{Si tu bisabuela viviera, ¿\textbf{le explicarías} qué es WhatsApp? ¿Cómo lo harías? ¿Con qué medio antiguo lo compararías para que entienda?}
\end{softbox}

\begin{softbox}{milcVerde}{milcGris}{Saber y saber hacer}
Al terminar la clase: (1) podrás \textbf{identificar} 8 medios de comunicación importantes en orden cronológico; (2) sabrás \textbf{analizar} qué cambió y qué se conservó entre dos medios; (3) habrás \textbf{creado} una línea de tiempo personal con dibujos y conexiones; (4) podrás \textbf{evaluar} con criterios cuál de los 8 medios trajo el cambio más profundo para tu comunidad.
\end{softbox}



\small
\begin{tabularx}{\textwidth}{>{\bfseries}p{3.0cm}Y}
\rowcolor{milcGradePrimary!12}Autor & Dr. Álvaro Cárdenas Orozco\\
DBA & Reconoce principios y conceptos propios de la tecnología, relacionando artefactos con su utilización segura (MEN, T\&I 6°)\\
Referentes & Saber ancestral del pregonero · Pensamiento computacional inicial · Cuidado de la palabra\\
Producto final & Una \textbf{línea de tiempo dibujada en hoja completa} del cuaderno, con \textbf{8 medios de comunicación} en orden cronológico (1850 hasta hoy). Cada medio lleva: año aproximado, dibujo pequeño, 1 frase de qué hacía, y \textbf{una flecha que conecta dos medios diciendo qué se conservó} (saludo, asunto, firma, audiencia, etc.). Más \textbf{una conclusión escrita} de 3-4 líneas: ¿qué cambió y qué sigue igual desde el pregonero hasta WhatsApp?\\
\end{tabularx}
\normalsize

\puente{Hoy haces 4 cosas: imaginas viajar en el tiempo, descubres los 8 medios principales, analizas qué se conservó, y dibujas tu línea de tiempo.}

\milcestacion{milcGradeDark}{Ruta MILC: así vamos a aprender}
\begin{tabularx}{\textwidth}{>{\centering\arraybackslash}X>{\centering\arraybackslash}X>{\centering\arraybackslash}X>{\centering\arraybackslash}X}
\phasecard{1}{milcVerde}{milcGris}{Escucha\\[-1mm]\small Observo, escucho y pregunto.} &
\phasecard{2}{milcTurquesa}{milcGris}{Entender\\[-1mm]\small Sistematizo y ordeno.} &
\phasecard{3}{milcMagenta}{milcGris}{Hacer\\[-1mm]\small Construyo y aplico.} &
\phasecard{4}{milcVino}{milcGris}{Cierre\\[-1mm]\small Evalúo y reflexiono.}
\end{tabularx}


\puente{Empezamos con un viaje mental al Cartago de tu bisabuela.}

\milcestacion{milcVerde}{Estación 1 de 3 · Escucha}

\begin{softbox}{milcVerde}{milcGris}{Escena para escuchar}
\textbf{Actividad 1 · IDENTIFICA --- Viaje al Cartago de 1920} (10 min · individual). Cierra los ojos un momento e imagina: estás en Cartago, año 1920. No hay celulares, no hay radio en tu casa, no hay televisión, no hay internet. Tu prima vive en Pereira. \textbf{¿Cómo le mandarías un mensaje?}. Después abre los ojos y \textbf{en el cuaderno} responde estas 3 preguntas con tus palabras: (1) ¿Qué medio usarías para mandarle un mensaje urgente? (2) ¿Cuánto tardaría en llegar? (3) ¿Cómo sabrías si lo recibió? Imagina con calma, no busques en internet.
\end{softbox}

{\sffamily\bfseries\fontsize{8}{10}\selectfont\textcolor{milcNegro!55}{MARCA CUANDO LO HAYAS HECHO}}
\milccheckrow{Cerré los ojos e imaginé el viaje al pasado.}
\milccheckrow{Respondí las 3 preguntas con mis palabras.}
\milccheckrow{No me ayudé de internet ni del compañero.}

\begin{infoband}{milcVerde}


\textbf{Lo que va al cuaderno (Actividad 1 · IDENTIFICA).} Título: «Actividad 1 --- Viaje a 1920». \textbf{Formato:} 3 preguntas con sus respuestas, en bloques cortos. \textbf{Extensión:} 3 párrafos de 2-3 líneas cada uno. \textbf{Sabes que terminaste cuando} te imaginaste tan bien la escena que pudiste responder sin dudar.
\end{infoband}

\puente{Después del viaje, ves los 8 medios principales con sus fechas y características.}

\milcestacion{milcTurquesa}{Estación 2 de 3 · Entender}

\begin{softbox}{milcTurquesa}{milcGris}{Lo que vas a entender}
Lo que imaginaste te conecta con el primer eslabón. Ahora vas a recorrer \textbf{8 medios de comunicación clave} desde 1850 hasta hoy. No tienes que memorizar las fechas exactas; lo importante es el orden y qué hacía cada uno.
\end{softbox}

\milcpaso{1}{milcTurquesa}{\textbf{1. Carta postal (siglo XIX, antes de 1900).} Mensajes en papel que viajaban en barco, en tren o en mula. Llegar a Pereira desde Cartago podía tardar 1 semana. \textbf{Lo que aportó:} la idea de \emph{destinatario, asunto, saludo y firma} --- las 5 partes del correo de hoy. \textbf{2. Telégrafo (1844 en EEUU, 1865 en Colombia).} Mensajes cortos que viajaban por cables eléctricos. Cada palabra costaba dinero, entonces se escribía con palabras pegadas: \emph{``LLEGOBIEN HERMANOENFERMO VENPRONTO''}. \textbf{Lo que aportó:} la idea de \emph{abreviar para enviar rápido} --- abuelo del SMS y del chat.}
\milcpaso{2}{milcTurquesa}{\textbf{3. Teléfono fijo (1876, Alexander Graham Bell).} Hablar con otra persona a distancia, en tiempo real, por voz. \textbf{Lo que aportó:} la idea de \emph{conversación a distancia} --- abuelo de las videollamadas. En Cartago llegó masivamente en los años 60. \textbf{4. Radio (años 1920 en Colombia).} Una sola persona habla y muchas personas escuchan al mismo tiempo. \textbf{Lo que aportó:} la idea de \emph{audiencia masiva} --- abuelo de TikTok y YouTube.}
\milcpaso{3}{milcTurquesa}{\textbf{5. Televisión (1954 en Colombia).} Como la radio, pero con imagen. La gente se sentaba a verla. \textbf{Lo que aportó:} la idea de \emph{ver + oír al mismo tiempo} --- abuelo del video. \textbf{6. Internet temprano (años 1990).} Computadores que se conectan unos con otros. En Cartago llegaron los café-internet en el 2000. \textbf{Lo que aportó:} la idea de \emph{cualquiera puede publicar} --- antes solo periódicos y TV publicaban.}
\milcpaso{4}{milcTurquesa}{\textbf{7. Celular y SMS (años 2000).} El teléfono móvil reemplazó al fijo. Mensajes de texto cortos. \textbf{Lo que aportó:} la idea de \emph{tener el medio en el bolsillo} --- el medio ya no está en una sala, está contigo. \textbf{8. Redes sociales y WhatsApp (2009-hoy).} Mezcla todo: voz como teléfono, imagen como TV, audiencia masiva como radio, texto como carta. \textbf{Lo que aportan:} la idea de \emph{participar como audiencia y como emisor a la vez} --- ya no solo escuchas, también respondes.}

\begin{drawbox}{milcTurquesa}{Línea de tiempo de los medios}
\textbf{1850} · Carta postal --- 1 semana de viaje. \\[1mm]
\textbf{1865} · Telégrafo en Colombia --- minutos, pero caro. \\[1mm]
\textbf{1876} · Teléfono fijo --- voz en tiempo real. \\[1mm]
\textbf{1920s} · Radio --- audiencia masiva. \\[1mm]
\textbf{1954} · Televisión --- voz e imagen. \\[1mm]
\textbf{1990s} · Internet --- todos publican. \\[1mm]
\textbf{2000s} · Celular y SMS --- medio en el bolsillo. \\[1mm]
\textbf{2009+} · Redes sociales y WhatsApp --- todo mezclado.
\end{drawbox}

\begin{softbox}{milcMostaza}{milcGris}{Errores comunes y su corrección}
\textbf{Confusiones frecuentes sobre la historia de los medios.} (1) \emph{``Internet inventó la comunicación masiva''} --- Falso, la radio y la TV ya la habían inventado. (2) \emph{``Antes de WhatsApp nadie sabía abreviar''} --- Falso, el telégrafo abreviaba hace 150 años. (3) \emph{``Las redes sociales son lo más nuevo que existe''} --- Cierto pero engañoso: las redes son \emph{una mezcla} de medios anteriores, no algo completamente nuevo. (4) \emph{``Las cartas son cosa muerta''} --- Falso, el correo electrónico es una carta y se sigue usando todos los días.
\end{softbox}

\begin{infoband}{milcTurquesa}


\textbf{Lo que va al cuaderno (Actividad 2 · EVALÚA).} Título: «Actividad 2 --- Cuál medio sirve mejor para qué». \textbf{Formato:} tabla de 3 columnas. Columna 1: característica. Columna 2: cómo era en la \emph{carta postal} de 1900. Columna 3: cómo es en \emph{WhatsApp} de hoy. Compara 5 características: velocidad, costo, audiencia, qué se conserva, qué se pierde. Al final escribe un \textbf{veredicto razonado de 3-4 líneas}: ¿cuál de los dos medios sirve mejor para mandarle a tu prima un mensaje urgente desde Cartago? ¿Y para celebrar los 90 años de tu abuela? Justifica cada elección con 2 razones. \textbf{Extensión:} 5 filas + veredicto. \textbf{Sabes que terminaste cuando} tu veredicto distingue \emph{rápido} de \emph{adecuado} --- no siempre lo más rápido es lo mejor.
\end{infoband}

\puente{Con los 8 medios claros, dibujas tu propia línea de tiempo con conexiones.}

\milcestacion{milcMagenta}{Estación 3 de 3 · Hacer}

\begin{softbox}{milcMagenta}{milcGris}{Cómo vas a construir tu producto}
\textbf{Actividad 3 · CREA --- Tu línea de tiempo personal} (25 min · individual). Vas a dibujar en \textbf{una hoja completa del cuaderno} una línea de tiempo con los 8 medios de comunicación, en orden cronológico. No es solo una lista: vas a agregar dibujos pequeños y \textbf{flechas que conectan medios} mostrando qué se conservó del uno al otro.
\end{softbox}

\milcpaso{1}{milcMagenta}{\textbf{Paso 1.} Pon la hoja \emph{horizontal} (apaisada). Dibuja \textbf{una línea horizontal larga} en el centro. En el extremo izquierdo escribe \textbf{1850}, en el extremo derecho escribe \textbf{2026}. Marca con líneas verticales pequeñas las fechas: 1850, 1865, 1876, 1920, 1954, 1990, 2000, 2010.}
\milcpaso{2}{milcMagenta}{\textbf{Paso 2.} En cada marca de fecha escribe \textbf{el nombre del medio} (Carta, Telégrafo, Teléfono, Radio, TV, Internet, SMS, WhatsApp). Alterna: unos arriba de la línea, otros abajo, para que no se amontonen.}
\milcpaso{3}{milcMagenta}{\textbf{Paso 3.} Al lado de cada nombre dibuja \textbf{un ícono pequeño} de ese medio: un sobre para la carta, una bocina para el teléfono, una antena para la radio, un televisor cuadrado para la TV, un computador para internet, un celular para SMS, un círculo verde con el rayo para WhatsApp. Dibujos simples, no tienen que estar perfectos.}
\milcpaso{4}{milcMagenta}{\textbf{Paso 4.} Bajo cada nombre escribe \textbf{una frase corta} (máximo 8 palabras) de qué hacía ese medio. Ejemplo: \emph{``Carta: viaja en mula, 1 semana''}, \emph{``Radio: una voz, miles oyendo''}.}
\milcpaso{5}{milcMagenta}{\textbf{Paso 5.} \textbf{Lo más importante:} dibuja \textbf{4 flechas curvas} que conecten un medio antiguo con uno nuevo, mostrando \textbf{qué se conservó}. Ejemplo: una flecha del telégrafo a SMS con la palabra \emph{``abreviar''}. Una flecha de la carta al correo electrónico con la palabra \emph{``saludo, asunto, firma''}. Tú decides las otras 2 conexiones.}
\milcpaso{6}{milcMagenta}{\textbf{Paso 6.} En la esquina inferior derecha escribe tu \textbf{conclusión personal} de 3-4 líneas. Empieza con: \emph{``En 176 años cambió...''} y termina con: \emph{``...pero sigue igual...''}. Ejemplo: \emph{``En 176 años cambió la velocidad, el alcance y el costo, pero sigue igual la necesidad de saludar, identificar al destinatario y firmar lo que escribimos.''}.}

\begin{drawbox}{milcMagenta}{Antes de mostrar la línea de tiempo}
\checkbox La línea ocupa una hoja completa horizontal. \\[1mm]
\checkbox Están los 8 medios con su fecha y nombre. \\[1mm]
\checkbox Cada medio tiene un ícono dibujado. \\[1mm]
\checkbox Cada medio tiene una frase corta de lo que hacía. \\[1mm]
\checkbox Hay al menos 4 flechas de conexión con su palabra. \\[1mm]
\checkbox Está la conclusión personal de 3-4 líneas firmada.
\end{drawbox}

\begin{softbox}{milcMagenta}{milcGris}{Plantilla de trabajo}
\textbf{Estructura del eje}: 1850 --- 1865 --- 1876 --- 1920s --- 1954 --- 1990s --- 2000s --- 2010+. \textbf{Cada medio}: fecha + nombre + ícono + frase corta. \textbf{Flechas}: conectan medio antiguo con moderno, llevan una palabra (lo que se conservó). \textbf{Conclusión}: 3-4 líneas en la esquina inferior derecha.
\end{softbox}

\begin{infoband}{milcMagenta}


\textbf{Lo que va al cuaderno (Actividad 3 · CREA).} Título: «Actividad 3 --- Mi línea de tiempo de los medios». \textbf{Formato:} hoja entera horizontal con línea de tiempo, dibujos pequeños, flechas y conclusión. \textbf{Extensión:} 1 hoja entera. \textbf{Sabes que terminaste cuando} podrías explicársela a tu bisabuela y ella entendería en 1 minuto cómo llegamos de su carta postal a tu WhatsApp.
\end{infoband}

\puente{El producto es esa línea de tiempo dibujada y la conclusión sobre qué se conservó.}

\milcestacion{milcMostaza}{Producto de la sesión}
\textbf{Línea de tiempo de los medios · 8 medios con conexiones · firmada}.
\begin{softbox}{milcMostaza}{milcGris}{Criterios de calidad}
(1) La línea ocupa una hoja entera horizontal. (2) Están los 8 medios en orden cronológico correcto. (3) Cada medio tiene fecha, nombre, ícono y frase corta. (4) Hay mínimo 4 flechas de conexión entre medios antiguos y modernos. (5) La conclusión final responde qué cambió y qué se conservó.
\end{softbox}

\puente{Antes de salir, miras tu línea y verificas que están los 8 medios y las flechas de conexión.}

\milcestacion{milcVino}{Evaluación liberadora · cinco dimensiones}

\begin{softbox}{milcVino}{milcGris}{Sentido de la evaluación}
Evaluar no es solo revisar una nota. Es reconocer qué aprendí, qué cambió en mí, qué emociones manejé, cómo me relaciono con la ciudad y la comunidad, y qué le dejaré a quien viene detrás.
\end{softbox}

\textbf{Mi reflexión liberadora en cinco dimensiones.} Completa cada frase iniciada:

\small
\begin{tabularx}{\textwidth}{>{\bfseries\RaggedRight\arraybackslash}p{3.4cm}Y>{\RaggedRight\arraybackslash}p{4.6cm}}
\rowcolor{milcVino}\color{white}Dimensión & \color{white}\bfseries Mi respuesta (frame starter) & \color{white}\bfseries Algo que descubrí\\
1. Personal & Lo que más cambió en mí fue \linead{6.5cm} & \linead{4.2cm}\\
2. Emocional & La emoción más fuerte fue \linead{2.5cm} y la regulé \linead{3cm} & \linead{4.2cm}\\
3. Ciudadana & En democracia esto importa porque \linead{6cm} & \linead{4.2cm}\\
4. Local (Cartago) & En mi barrio sería útil para \linead{6.5cm} & \linead{4.2cm}\\
5. Intergeneracional & A mi abuela/abuelo le diría \linead{6.5cm} & \linead{4.2cm}\\
\end{tabularx}
\normalsize

\begin{infoband}{milcVino}
\textbf{Para la próxima sustentación.} Vuelve a esta tabla en seis meses. Verás que las respuestas cambian — no porque lo aprendido cambie, sino porque tú cambias. La evaluación liberadora es una conversación contigo a través del tiempo.
\end{infoband}


\puente{Cerramos con 3 voces que te ayudan a entender por qué saber historia te protege en el presente.}

\milcestacion{milcGradeDark}{Triángulo de pensamiento · cierre filosófico}

\begin{softbox}{milcVino}{milcGris}{Dussel · lente del nosotros}
\textit{El que no sabe de dónde viene una herramienta, fácil cae en pensar que la herramienta lo gobierna a él, y no al revés.}

{\footnotesize\itshape\color{milcNegro!70}--- formulación de esta guía, al modo de Enrique Dussel. No es una cita textual.}

Cuando \textbf{sabes que WhatsApp viene de la carta postal}, dejas de tratarlo como magia. Es un medio. Tiene reglas viejas (el saludo, la firma, el destinatario). Tú lo manejas, no él a ti. \textbf{La historia te devuelve poder sobre la herramienta}: te muestra que los humanos llevan 200 años pasándose mensajes y que tú no eres la primera generación con ese desafío.

\textbf{¿Uso las herramientas digitales con criterio, o dejo que ellas me usen a mí porque no sé de dónde vienen?}
\end{softbox}
\linead{\linewidth}\\[1mm]
\linead{\linewidth}

\begin{softbox}{milcMostaza}{milcGris}{Estoicismo · Marco Aurelio · cuidado interior}
\textit{Las formas cambian rápido. Lo que cuenta --- el saludo, la honestidad, el respeto --- es lo mismo desde hace milenios.}

{\footnotesize\itshape\color{milcNegro!70}--- formulación de esta guía, al modo de Marco Aurelio. No es una cita textual.}

Mira tu línea de tiempo: \textbf{las formas cambiaron} --- mula, cable, ondas, satélite, fibra óptica --- pero \textbf{lo importante sigue igual}. Se saluda, se identifica el destinatario, se firma, se respeta al otro. Si dominas eso, dominas cualquier medio nuevo que aparezca dentro de 10 años. La forma cambiará otra vez. Lo bueno --- el respeto --- no.

\textbf{¿Estoy aprendiendo lo que dura, o solo lo que está de moda esta semana?}
\end{softbox}
\linead{\linewidth}\\[1mm]
\linead{\linewidth}

\begin{softbox}{milcTurquesa}{milcGris}{Floridi · lente de la infoesfera}
\textit{Cada generación cree que su medio es el más importante. Después llega otro medio y se vuelven el medio anterior. La sabiduría es ver venir el cambio.}

{\footnotesize\itshape\color{milcNegro!70}--- formulación de esta guía, al modo de Luciano Floridi. No es una cita textual.}

Tu bisabuela vivió un mundo \textbf{sin radio}. Tu abuela vivió un mundo \textbf{sin TV}. Tu mamá vivió un mundo \textbf{sin WhatsApp}. Tú vas a vivir un mundo \textbf{sin algún medio que apenas se va a inventar}. La pregunta no es \emph{``qué medio voy a usar siempre''}: la pregunta es \emph{``cómo aprendo a aprender medios nuevos sin enredarme''}. La línea de tiempo te enseña eso.

\textbf{¿Estoy aprendiendo a \textbf{cambiarme} de medio cuando llegue uno nuevo, o me voy a quedar pegado al que conozco hoy?}
\end{softbox}
\linead{\linewidth}\\[1mm]
\linead{\linewidth}

\begin{infoband}{milcNegro}
\textbf{Nota para el docente.} Las tres formulaciones son del autor de la guía, no citas textuales de los pensadores: recogen su lente para pensar el tema, no sus palabras. Si vas a citarlos en clase, remite a la obra original.
\end{infoband}

\milcestacion{milcVino}{Mi autoevaluación · marca una casilla por fila}

{\small
\setlength{\extrarowheight}{2.2pt}
\begin{tabularx}{\textwidth}{>{\RaggedRight\arraybackslash\bfseries}p{3.0cm}YYY}
\rowcolor{milcVino}\color{white}\bfseries Criterio & \color{white}\bfseries Lo logré & \color{white}\bfseries En proceso & \color{white}\bfseries Necesito apoyo\\[.6mm]
Comprensión del tema & \milcopcion{Explico las ideas con mis palabras.} & \milcopcion{Reconozco las ideas, falta precisión.} & \milcopcion{Necesito repasar lo básico.}\\[.8mm]
\rowcolor{milcGris}Pensamiento computacional & \milcopcion{Usé los pilares al armar mi producto.} & \milcopcion{Usé algunos pilares.} & \milcopcion{Necesito releer la estación 2.}\\[.8mm]
Aplicación y producto & \milcopcion{Mi producto cumple los criterios.} & \milcopcion{Está armado, con ajustes pendientes.} & \milcopcion{Aún está en construcción.}\\[.8mm]
\rowcolor{milcGris}Reflexión 5 dimensiones & \milcopcion{Respondí cada una con honestidad.} & \milcopcion{Respondí algunas con detalle.} & \milcopcion{Mis respuestas son muy generales.}\\[.8mm]
Triángulo de pensamiento & \milcopcion{Conecté las 3 voces con mi trabajo.} & \milcopcion{Conecté al menos una.} & \milcopcion{No las relaciono con mi proceso.}\\[.8mm]
\end{tabularx}}

\vspace{3mm}
\textbf{Hoy entendí que:}\\[1mm]
\linead{\linewidth}\\[1mm]
\linead{\linewidth}

\textbf{Mi compromiso para la próxima sesión es:}\\[1mm]
\linead{\linewidth}\\[1mm]
\linead{\linewidth}

\begin{softbox}{milcMostaza}{milcGris}{Cita de despedida · Marco Aurelio}
\textit{``El obstáculo en el camino se vuelve el camino. Lo que estorba al camino se vuelve camino.''}\\[1mm]
Cada error de hoy, bien observado, se convierte en pieza del camino que pisarás mañana.
\end{softbox}

\small
\begin{tabularx}{\textwidth}{>{\bfseries}p{3.6cm}Y}
\rowcolor{milcGradeDark!12}Nombre completo: & \linead{10cm}\\
Fecha: & \linead{4cm} \quad Curso/grupo: \linead{4cm}\\
Mi firma: & \linead{9cm}\\
\end{tabularx}
\normalsize

\begin{infoband}{milcGradeDark}
\textbf{Para la próxima sesión.} Guarda tu línea de tiempo en el cuaderno. Esta semana, observa: ¿cuántos medios distintos usaste en un día? Cuéntalos. La próxima clase vas a aprender sobre tu \textbf{huella digital}: el rastro de cosas que dejas cuando navegas, posteas y juegas. Hoy viste la historia. La próxima clase vas a ver el presente: lo que va quedando con tu nombre.
\end{infoband}

\end{document}
